% helden/gorbas.tex - weltlicher Held ohne Magie und ohne Liturgien.
% Probe auf die Automatik -OhneLeere: Bogen 4 und 5 muessen wegfallen.
% helden/dorle.tex — Werte eines Helden. Nur Werte, keine Koordinaten.
% Fehlt ein Wert, bleibt das Feld leer. Kein Abbruch.
% Die Feldnamen sind die der Feldtabelle und in beiden Quellfassungen gleich.

\wert{s1_name}{Gorbas Sohn des Gorm}
\wert{s1_geschlecht}{weiblich}
\wert{s1_kultur}{Mittelländisches Dorf}
% s1_profession absichtlich nicht gesetzt — Probe auf den Leerfall

% Abgeleitete Werte. Der Bogen hat die Formeln aufgedruckt; hier stehen nur
% Zahlen, es wird nichts nachgerechnet.
\wert{s1_lep_wert}{32}
\wert{s1_lep_bonus}{2}
\wert{s1_lep_max}{34}
\wert{s1_sk_wert}{1}
\wert{s1_sk_max}{1}
\wert{s1_zk_wert}{2}
\wert{s1_zk_max}{2}
\wert{s1_aw_wert}{5}
\wert{s1_aw_max}{5}
\wert{s1_schips_wert}{3}
\wert{s1_schips_max}{3}
% AsP und KaP absichtlich leer — Dorle ist weder Zauberin noch Geweihte

% Restliche Angaben Seite 1
\wert{s1_spezies}{Mensch}
\wert{s1_sozialstatus}{unfrei}
\wert{s1_heimatort}{Fuchsgrund}
\wert{s1_alter}{34}
\wert{s1_haarfarbe}{graubraun}
\wert{s1_augenfarbe}{grau}
\wert{s1_groesse}{1,64 m}
\wert{s1_gewicht}{58 kg}
% s1_familie, s1_charakter_* absichtlich leer

% --- Bogen 3: Kampfwerte ---------------------------------------------------
\wert{s3_gs}{8}
\wert{s3_le}{32}
\wert{s3_aw}{5}
\wert{s3_ini}{12}
\wert{s3_sk}{1}
\wert{s3_zk}{2}

\wert{s3_kt_dolche_ktw}{10}
\wert{s3_kt_dolche_atfk}{8}
\wert{s3_kt_dolche_pa}{5}
\wert{s3_kt_raufen_ktw}{8}
\wert{s3_kt_raufen_atfk}{7}
\wert{s3_kt_raufen_pa}{4}
\wert{s3_kt_armbrueste_ktw}{6}
\wert{s3_kt_armbrueste_atfk}{6}
\wert{s3_kt_schwerter_ktw}{12}
\wert{s3_kt_schwerter_atfk}{9}
\wert{s3_kt_schwerter_pa}{6}
\wert{s3_kt_wurfwaffen_ktw}{7}
\wert{s3_kt_wurfwaffen_atfk}{7}

\wert{s3_ksf_1}{Wuchtschlag I, Finte I}
\wert{s3_ksf_2}{Klingensturm}

\wert{s3_lep_max}{34}
\wert{s3_lep_aktuell}{34}
\wert{s3_lep_schwelle_1}{25}
\wert{s3_lep_schwelle_2}{17}
\wert{s3_lep_schwelle_3}{8}
\wert{s3_lep_schwelle_4}{5}

% --- Bogen 3: Tabellen (Testwerte zur Sichtprobe der Spalten) -------------
\wert{s3_nk_1_waffe}{Dolch}
\wert{s3_nk_1_technik}{Dolche}
\wert{s3_nk_1_schaden}{KK 14/3}
\wert{s3_nk_1_tp_basis}{4}
\wert{s3_nk_1_tp_gesamt}{5}
\wert{s3_nk_1_mod}{0/0}
\wert{s3_nk_1_reichweite}{kurz}
\wert{s3_nk_1_at}{9}
\wert{s3_nk_1_pa}{5}
\wert{s3_nk_1_gewicht}{1 Unz}

\wert{s3_fk_1_waffe}{Schwere Armbrust}
\wert{s3_fk_1_technik}{Armbrueste}
\wert{s3_fk_1_ladezeit}{16 Akt}
\wert{s3_fk_1_tp}{13}
\wert{s3_fk_1_reichweite}{25/50/100}
\wert{s3_fk_1_fernkampf}{6}
\wert{s3_fk_1_munition}{12 Bolzen}
\wert{s3_fk_1_gewicht}{5 Stn}

\wert{s3_ru_1_name}{Lederharnisch}
\wert{s3_ru_1_rs}{3}
\wert{s3_ru_1_be}{1}
\wert{s3_ru_1_abzuege}{-}
\wert{s3_ru_1_gewicht}{3 Stn}
\wert{s3_ru_1_wo}{Torso}

\wert{s3_sch_1_name}{Holzschild}
\wert{s3_sch_1_struktur}{4}
\wert{s3_sch_1_mod}{0/+1}
\wert{s3_sch_1_gewicht}{2 Stn}

\wert{s3_zu_1_stufe1}{X}
\wert{s3_zu_1_stufe2}{X}
\wert{s3_zu_1_stufe3}{X}
\wert{s3_zu_1_stufe4}{X}
\wert{s3_zu_11_stufe1}{X}

% --- Bogen 1: Eigenschaften -----------------------------------------------
\wert{s1_mu}{12}
\wert{s1_kl}{11}
\wert{s1_in}{13}
\wert{s1_ch}{10}
\wert{s1_ff}{12}
\wert{s1_ge}{13}
\wert{s1_ko}{12}
\wert{s1_kk}{11}

% --- Bogen 2: Talente (Stichproben in allen fünf Gruppen) ----------------
\wert{s2_belastung}{1}
\wert{s2_mu}{12} \wert{s2_kl}{11} \wert{s2_in}{13} \wert{s2_ch}{10}
\wert{s2_ff}{12} \wert{s2_ge}{13} \wert{s2_ko}{12} \wert{s2_kk}{11}
\wert{s2_fliegen_fw}{0}
\wert{s2_klettern_fw}{6}
\wert{s2_klettern_rpr}{-}
\wert{s2_klettern_anm}{Fels}
\wert{s2_wildnisleben_fw}{9}
\wert{s2_wildnisleben_rpr}{3}
\wert{s2_wildnisleben_anm}{Wald}
\wert{s2_willenskraft_fw}{4}
\wert{s2_brettspiel_fw}{3}
\wert{s2_sternkunde_fw}{2}
\wert{s2_alchimie_fw}{1}
\wert{s2_stoffbearbeitung_fw}{7}
\wert{s2_stoffbearbeitung_rpr}{2}
\wert{s2_stoffbearbeitung_anm}{Flicken}
\wert{s2_mod_mu_0}{X}
\wert{s2_mod_kk_p2}{X}
\wert{s2_sprachen_1}{Garethi (Muttersprache)}
\wert{s2_sprachen_2}{Zwergisch (Grundkenntnisse)}
\wert{s2_schriften_1}{Kusliker Zeichen}

% --- Bogen 6: Ausrüstung --------------------------------------------------
\wert{s6_gg_l_1_gegenstand}{Wanderstab}
\wert{s6_gg_l_1_gewicht}{2 Stn}
\wert{s6_gg_l_1_wo}{Hand}
\wert{s6_gg_l_30_gegenstand}{letzte Zeile links}
\wert{s6_gg_r_1_gegenstand}{Wollmantel}
\wert{s6_gg_r_1_gewicht}{3 Stn}
\wert{s6_gg_r_1_wo}{getragen}
\wert{s6_gg_r_30_gegenstand}{letzte Zeile rechts}
\wert{s6_gesamtgewicht_l}{14 Stn}
\wert{s6_gesamtgewicht_r}{9 Stn}
\wert{s6_tragkraft_l}{22}
\wert{s6_tragkraft_r}{22}
\wert{s6_tk_basis}{22}
\wert{s6_tk_4}{26}
\wert{s6_tk_8}{30}
\wert{s6_tk_12}{34}
\wert{s6_tk_16}{38}
\wert{s6_dukaten}{2}
\wert{s6_silbertaler}{14}
\wert{s6_heller}{7}
\wert{s6_kreuzer}{3}
\wert{s6_tier_name}{Grauschnauze}
\wert{s6_tier_lo}{2}
\wert{s6_tier_ap}{80}
\wert{s6_tier_typus}{Haushund}
\wert{s6_tier_lep}{18}
\wert{s6_tier_mu}{13}
\wert{s6_tier_kl}{6}
\wert{s6_tier_in}{14}
\wert{s6_tier_ch}{11}
\wert{s6_tier_ff}{8}
\wert{s6_tier_ge}{16}
\wert{s6_tier_ko}{13}
\wert{s6_tier_kk}{11}
\wert{s6_tier_sk}{2}
\wert{s6_tier_zk}{3}
\wert{s6_tier_rs}{0}
\wert{s6_tier_ini}{15}
\wert{s6_tier_gs}{12}
\wert{s6_tier_angriff}{Biss}
\wert{s6_tier_at}{11}
\wert{s6_tier_vw}{5}
\wert{s6_tier_tp}{1W6}
\wert{s6_tier_rw}{kurz}
\wert{s6_tier_aktionen}{1 Angriff}
\wert{s6_tier_sf}{Gutes Gehoer}

% Eigenschaftsleiste auf Bogen 4 und 5 (zur Prüfung der Zellposition)
\wert{s4_mu}{12} \wert{s4_kl}{11} \wert{s4_in}{13} \wert{s4_ch}{10}
\wert{s4_ff}{12} \wert{s4_ge}{13} \wert{s4_ko}{12} \wert{s4_kk}{11}
\wert{s5_mu}{12} \wert{s5_kl}{11} \wert{s5_in}{13} \wert{s5_ch}{10}
\wert{s5_ff}{12} \wert{s5_ge}{13} \wert{s5_ko}{12} \wert{s5_kk}{11}

% Bogen 1: Textblöcke
\wert{s1_vorteile_1}{Zäher Hund, Ortskenntnis (Fuchsgrund)}
\wert{s1_nachteile_1}{Streitsucht, Neugier}
\wert{s1_sf_1}{Ortskenntnis, Wildnisleben-Anwendung}
\wert{s1_schips_aktuell}{3}
\wert{s1_erfahrungsgrad}{Erfahren}
\wert{s1_ap_gesamt}{1100}
\wert{s1_ap_gesammelt}{1100}
\wert{s1_ap_ausgegeben}{1050}

% Heldenbild. Pfad relativ zum Projektverzeichnis. Fehlt die Datei, bleibt
% der Rahmen leer und der Lauf gibt nur einen Hinweis ins Log.
\wert{s1_bild}{helden/bilder/pruefbild.png}
